% Options for packages loaded elsewhere
\PassOptionsToPackage{unicode}{hyperref}
\PassOptionsToPackage{hyphens}{url}
\documentclass[
  10pt,
  a4paper,
]{article}
\usepackage{xcolor}
\usepackage[margin=22mm]{geometry}
\usepackage{amsmath,amssymb}
\setcounter{secnumdepth}{-\maxdimen} % remove section numbering
\usepackage{iftex}
\ifPDFTeX
  \usepackage[T1]{fontenc}
  \usepackage[utf8]{inputenc}
  \usepackage{textcomp} % provide euro and other symbols
\else % if luatex or xetex
  \usepackage{unicode-math} % this also loads fontspec
  \defaultfontfeatures{Scale=MatchLowercase}
  \defaultfontfeatures[\rmfamily]{Ligatures=TeX,Scale=1}
\fi
\usepackage{lmodern}
\ifPDFTeX\else
  % xetex/luatex font selection
\fi
% Use upquote if available, for straight quotes in verbatim environments
\IfFileExists{upquote.sty}{\usepackage{upquote}}{}
\IfFileExists{microtype.sty}{% use microtype if available
  \usepackage[]{microtype}
  \UseMicrotypeSet[protrusion]{basicmath} % disable protrusion for tt fonts
}{}
\makeatletter
\@ifundefined{KOMAClassName}{% if non-KOMA class
  \IfFileExists{parskip.sty}{%
    \usepackage{parskip}
  }{% else
    \setlength{\parindent}{0pt}
    \setlength{\parskip}{6pt plus 2pt minus 1pt}}
}{% if KOMA class
  \KOMAoptions{parskip=half}}
\makeatother
\setlength{\emergencystretch}{3em} % prevent overfull lines
\providecommand{\tightlist}{%
  \setlength{\itemsep}{0pt}\setlength{\parskip}{0pt}}
\usepackage{mathrsfs}
\newcommand{\Log}{\operatorname{Log}}
\usepackage{mathtools}
\usepackage{xurl}
\emergencystretch=3em
\setcounter{tocdepth}{1}
\newcommand{\Beta}{\mathrm{B}}
\DeclareUnicodeCharacter{00B2}{\ensuremath{^2}}
\DeclareUnicodeCharacter{00B1}{\ensuremath{\pm}}
\DeclareUnicodeCharacter{00D7}{\ensuremath{\times}}
\DeclareUnicodeCharacter{2192}{\ensuremath{\to}}
\DeclareUnicodeCharacter{21D2}{\ensuremath{\Rightarrow}}
\DeclareUnicodeCharacter{21A6}{\ensuremath{\mapsto}}
\DeclareUnicodeCharacter{2208}{\ensuremath{\in}}
\DeclareUnicodeCharacter{2209}{\ensuremath{\notin}}
\DeclareUnicodeCharacter{2212}{\ensuremath{-}}
\DeclareUnicodeCharacter{2227}{\ensuremath{\wedge}}
\DeclareUnicodeCharacter{2228}{\ensuremath{\vee}}
\DeclareUnicodeCharacter{2260}{\ensuremath{\ne}}
\DeclareUnicodeCharacter{2264}{\ensuremath{\le}}
\DeclareUnicodeCharacter{2265}{\ensuremath{\ge}}
\DeclareUnicodeCharacter{2297}{\ensuremath{\otimes}}
\DeclareUnicodeCharacter{00B3}{\ensuremath{^3}}
\DeclareUnicodeCharacter{2074}{\ensuremath{^4}}
\DeclareUnicodeCharacter{2076}{\ensuremath{^6}}
\DeclareUnicodeCharacter{03BB}{\ensuremath{\lambda}}
\DeclareUnicodeCharacter{03BC}{\ensuremath{\mu}}
\DeclareUnicodeCharacter{03B5}{\ensuremath{\epsilon}}
\DeclareUnicodeCharacter{03C4}{\ensuremath{\tau}}
\DeclareUnicodeCharacter{03A6}{\ensuremath{\Phi}}
\DeclareUnicodeCharacter{03B1}{\ensuremath{\alpha}}
\DeclareUnicodeCharacter{03C3}{\ensuremath{\sigma}}
\DeclareUnicodeCharacter{03B4}{\ensuremath{\delta}}
\DeclareUnicodeCharacter{03C0}{\ensuremath{\pi}}
\DeclareUnicodeCharacter{039B}{\ensuremath{\Lambda}}
\DeclareUnicodeCharacter{2202}{\ensuremath{\partial}}
\DeclareUnicodeCharacter{0393}{\ensuremath{\Gamma}}
\DeclareUnicodeCharacter{2205}{\ensuremath{\varnothing}}
\DeclareUnicodeCharacter{221E}{\ensuremath{\infty}}
\DeclareUnicodeCharacter{2200}{\ensuremath{\forall}}
\DeclareUnicodeCharacter{00B9}{\ensuremath{^1}}
\DeclareUnicodeCharacter{00B0}{\ensuremath{{}^\circ}}
\DeclareUnicodeCharacter{211D}{\ensuremath{\mathbb{R}}}
\usepackage{bookmark}
\IfFileExists{xurl.sty}{\usepackage{xurl}}{} % add URL line breaks if available
\urlstyle{same}
\hypersetup{
  hidelinks,
  pdfcreator={LaTeX via pandoc}}

\author{}
\date{}

\begin{document}

{
\setcounter{tocdepth}{3}
\tableofcontents
}
\section{Full Weil coercivity across the first prime
window}\label{full-weil-coercivity-across-the-first-prime-window}

\textbf{Original-zeta receiving calculation.} The working meromorphic
function is the original Riemann zeta, with its full Gamma/endpoints
multiplier and its full trivial-zero and pole divisor retained.
\href{FAITHFUL_THETA_COMPLETION_RETURN.md}{TF1--42} proves the original
labelled theta inverse and the actual heat-image defect.
\href{FAITHFUL_UNCOMPLETED_ZETA_HEAT.md}{UZ1--53} proves every
exceptional-point fibre, jet, original-zeta heat term and reflection
orientation.
\href{ORIGINAL_ZETA_COMPACT_WEIL_RECONSTRUCTION.md}{OZC1--48} rederives
the compact contour and interval operator directly from zeta, including
the left-cutoff Gamma boundary, and specifies exactly which compensated
pairing the bounds concern.
\href{ORIGINAL_ZETA_HEAT_CONTACT_RECONSTRUCTION.md}{OZH1--49} rederives
the actual meromorphic heat, rational signed trace, compact-test domain,
contact drift and full causal arithmetic variation. The raw full divisor
and the compensated Weil receiver are linked by their displayed
correction, not identified. In particular the fixed-test trivial-zero
sum converges exactly at translations at least 1/32 and equals the
earlier R term; at zero translation it requires the proved cutoff
compensation. \href{ORIGINAL_ZETA_CAUCHY_RECONSTRUCTION.md}{OZK1--38}
reconstructs the original signed Cauchy trace, its full Gamma
correction, resonant finite parts, every finite matrix and its index,
and the complete local heat jets. Its invertible map retains the raw
trace and correction separately. These complete receiving proofs govern
the interpretation of the retained auxiliary calculations below.

The full reflected Weil form has an explicit positive lower bound on
every compact smooth test satisfying its two global endpoint equations
and supported in an interval of length at most \(19/25\). The
prime-\(2\) interaction is included. Its two overlap regions are
controlled by the exact archimedean boundary potential; the global
moment equations then control the two lowest Legendre modes.

Precisely, put \[
T_*=\frac{19}{25},\qquad c_*=\frac{11509}{600000}.
\tag{FC1}
\] For every nonzero complex-valued test in \[
\mathcal V_{T_*}=
\left\{F\in C_c^\infty(\mathbb R;\mathbb C):
\int F(t)e^{t/2}\,dt=\int F(t)e^{-t/2}\,dt=0,\
\operatorname{diam}(\operatorname{supp}F)\le T_*\right\},
\] the result is \[
\boxed{Q(F)>c_*\|F\|_2^2>\frac1{64}\|F\|_2^2.}
\tag{FC2}
\] For the zero function all terms are zero; the corresponding
non-strict inequality holds on the whole space.

This extends the supplied companion manuscript's \(3/4\)-support
calculation N48--N52. That incoming result and proof strategy are
retained as the source of the starting estimate; the larger support
width, its explicit constant, and the application of that width to the
entire first-prime window are the continuation here. Every estimate
needed for (FC2) is proved below. No numerical location of a zeta zero
is used, and no assertion of global Weil positivity or RH follows from
this local theorem.

\subsection{1. Original form, involution, and all endpoint
coordinates}\label{original-form-involution-and-all-endpoint-coordinates}

Use exactly \[
M_F(s)=\int_{\mathbb R}F(t)e^{-(s-1/2)t}\,dt,\qquad
F^\#(t)=\overline{F(-t)},\qquad
\tau_uF(t)=F(t-u),\qquad
\langle F,G\rangle=\int\overline F G.
\tag{FC3}
\] For \(h=F^\#*F\), \[
h(u)=\int\overline{F(t-u)}F(t)\,dt,\qquad
M_h(s)=\overline{M_F(1-\overline s)}M_F(s).
\] The full explicit formula, with every nontrivial zero and
multiplicity, is \[
\begin{aligned}
Q(F)&=\sum_\rho m_\rho
\overline{M_F(1-\overline\rho)}M_F(\rho)\\
&=M_h(0)+M_h(1)+A_\infty(h)-P_{\rm fin}(h),\\
P_{\rm fin}(h)&=
2\sum_{n\ge2}\frac{\Lambda(n)}{\sqrt n}
\Re\langle F,\tau_{\log n}F\rangle.
\end{aligned}
\tag{FC4}
\] Here \(\Lambda\) is the von Mangoldt function, with
\(\Lambda(p^j)=\log p\) and zero otherwise. Compact support makes the
prime sum finite. The original critical-strip zero count and rapid
vertical decrease of the compact smooth transforms prove absolute
convergence of the zero sum. The complete derivation, including the full
supported endpoints, is
\href{https://github.com/KokunoYumeto/zeta-function-research-reader/blob/5a89872598df0902b7c1393cf8e4692ca3010a95/workbenches/splitzero-tandem/branches/tau-zero-prime-positivity/SUPPORTED_ZERO_PRIME_WEIL_DERIVATION.md}{SZW19--SZW38}.

The two original endpoint equations imply \[
M_h(0)=\overline{M_F(1)}M_F(0)=0,\qquad
M_h(1)=\overline{M_F(0)}M_F(1)=0.
\tag{FC5}
\] They are global equations on \(F\). No separate moment equations will
be imposed on pieces of its support.

For comparison with the incoming plus-exponent convention, the exact map
is reflection: \[
(RF)(t)=F(-t),\qquad R^2=1,\qquad
M_F^+(s)=M_F(1-s)=M_{RF}(s).
\tag{FC6}
\] It is unitary on \(L^2\), preserves support diameter, and exchanges
the moments. It preserves \(Q\), as follows by relabelling the
absolutely convergent zero sum by \(\rho\mapsto1-\rho\). Thus the
original convention and the supplied convention are related by a proved
map, with no change to the Weil involution.

\subsection{2. The complete archimedean
expression}\label{the-complete-archimedean-expression}

With \(N=\|F\|_2^2\), the original Gamma contribution is \[
A_\infty(h)=-(\gamma_E+\log\pi)N+
\int_0^\infty
\frac{e^{-x}N-\tfrac12e^{-x/4}\{h(x/2)+h(-x/2)\}}
{1-e^{-x}}\,dx.
\tag{FC7}
\] This is HA13, also PS3, derived from the exact multiplier
\(\Re\psi(1/4+iy/2)-\log\pi\) in SZW24. The subtracted numerator makes
the integral converge at zero.

Put \[
k(u)=\frac{e^{-u/2}}{1-e^{-2u}}=\frac{e^{u/2}}{2\sinh u},
\qquad
\kappa_\infty=\log(8\pi)+\gamma_E+\frac\pi2.
\] Changing \(x=2u\) in (FC7) and using \(\|F-\tau_uF\|^2=2N-2\Re h(u)\)
gives \[
A_\infty(h)=
\int_0^\infty k(u)\|F-\tau_uF\|_2^2\,du-\kappa_\infty N.
\tag{FC8}
\] Here is an explicit scalar verification. First (FC7) is the negative
of \[
\{\log(4\pi)+\gamma_E\}N+
2\int_0^\infty\{\Re h(u)-e^{-u/2}N\}k(u)\,du.
\] The sum of those two expressions is \[
\log4\,N+
2N\int_0^\infty
\frac{e^{-2u}-e^{-u}}{1-e^{-2u}}\,du=0.
\] Next, with \(q=e^{-u/2}\), \[
2\int_0^\infty(1-e^{-u/2})k(u)\,du
=4\int_0^1\frac{dq}{(1+q)(1+q^2)}
=\log2+\frac\pi2.
\tag{FC9}
\] The partial fractions are \(\{2(1+q)\}^{-1}+(1-q)/\{2(1+q^2)\}\).
These identities prove (FC8), including the entire constant.

Equations (FC4)--(FC9) give the original full form on the two-moment
space: \[
Q(F)=\int_0^\infty k(u)\|F-\tau_uF\|_2^2\,du
-\kappa_\infty\|F\|_2^2
-2\sum_{n\ge2}\frac{\Lambda(n)}{\sqrt n}
\Re\langle F,\tau_{\log n}F\rangle.
\tag{FC10}
\] Near zero, \(\|F-\tau_uF\|_2\le u\|F'\|_2\); at infinity \(k\) decays
exponentially. This verifies the convergence used above. In particular,
the endpoint/Gamma term has not been replaced by a constant-only
approximation.

\subsection{3. Exact interval transport and exterior
potential}\label{exact-interval-transport-and-exterior-potential}

For a containing interval \(I=[c-T/2,c+T/2]\), use the unitary map \[
f(x)=(U_{T,c}F)(x)=\sqrt{\frac T2}\,F(c+Tx/2),
\qquad -1\le x\le1,
\] \[
F(t)=\sqrt{\frac2T}\,f(2(t-c)/T),\qquad
\|f\|_2=\|F\|_2.
\tag{FC11}
\] Translation preserves the full form:
\(M_{\tau_cF}(s)=e^{-c(s-1/2)}M_F(s)\), and the two factors in the
reflected pairing cancel. The transformed moments are exactly \[
\int_{-1}^1f(x)e^{ax}\,dx=
\int_{-1}^1f(x)e^{-ax}\,dx=0,\qquad a=T/4.
\tag{FC12}
\] The original scalar factors from the centre and Jacobian are nonzero.
The displayed equations are therefore equivalent to the original two
moments.

Let \[
\begin{aligned}
(\mathcal Lf)(x)&=\frac12\int_{-1}^1
\frac{f(x)-f(y)}{|x-y|}\,dy,\\
v(x)&=-\frac12\log(1-x^2),\qquad
r(s)=k(s)-\frac1{2s},\\
(\mathcal R_Tf)(x)&=\frac T2\int_{-1}^1
r(T|x-y|/2)f(y)\,dy.
\end{aligned}
\tag{FC13}
\] The exact archimedean operator transported by (FC11) is \[
\boxed{
A_{\infty,T}=\mathcal L+v-\log(\pi T)-\gamma_E-\mathcal R_T.
}
\tag{FC14}
\] To prove it, write the translation energy as
\(\frac12\int_{\mathbb R^2}k(|t-u|)|F(t)-F(u)|^2\,dt\,du\). Interior
pairs give the difference kernel after transport. Exterior pairs give
the diagonal potential \[
J(T(1+x)/2)+J(T(1-x)/2),\qquad
J(s)=\int_s^\infty k(u)\,du.
\] Direct differentiation and the limit at zero prove \[
J(s)=\operatorname{arctanh}(e^{-s/2})
+\arctan(e^{-s/2})
=-\frac12\log s+\log2+\frac\pi4-\int_0^s r(u)\,du.
\tag{FC15}
\] The interior diagonal part of \(r\) cancels the two integrals in
(FC15). Its off-diagonal part is \(-\mathcal R_T\). The remaining
exterior constant is \(-\log(T/2)+2\log2+\pi/2\). Subtracting
\(\kappa_\infty\) leaves \(-\log(\pi T)-\gamma_E\); the remaining
variable term is \(v(x)\). This proves (FC14), retaining the complete
exterior contribution.

\subsection{\texorpdfstring{4. Control the actual prime-\(2\) term for a
range of
widths}{4. Control the actual prime-2 term for a range of widths}}\label{control-the-actual-prime-2-term-for-a-range-of-widths}

For this section let \[
\frac34\le T\le\frac{19}{25},\qquad
\beta=\frac{\log2}{\sqrt2},\qquad
\delta=\frac{2\log2}{T}.
\tag{FC16}
\] One has \(\log2<T<\log3\) and \(1<\delta<2\). Thus only the prime
power \(2\) occurs in (FC10). Its full transported value is \[
\langle f,B_2f\rangle=
2\beta\Re\int_{-1+\delta}^1\overline{f(x)}f(x-\delta)\,dx
\le\beta\int_{|x|\ge\delta-1}|f(x)|^2\,dx.
\tag{FC17}
\] The inequality is \(2\Re(\overline zw)\le|z|^2+|w|^2\). The overlap
strips \([-1,1-\delta]\) and \([-1+\delta,1]\) are disjoint; changing
variables in the second squared term gives exactly the displayed
integral. No phase or part of the prime interaction is omitted.

The convergent positive series \[
\log2=2\sum_{j=0}^\infty\frac1{(2j+1)3^{2j+1}}
\tag{FC18}
\] follows by integrating the geometric series for \(1/(1-x^2)\) from
zero to \(1/3\). Its first two terms give \(\log2>56/81\). Therefore
uniformly over (FC16), \[
\delta-1>\frac{1261}{1539},\qquad
\left(\frac{1261}{1539}\right)^2-\frac7{11}
=\frac{911684}{26053731}>0.
\] The bound \(e<11/4\) gives \(1-e^{-1}<7/11\), so \(v(x)>1/2\)
throughout both overlap strips. Outside them \(v(x)\ge0\).

Also \(e^{7/10}>12013/6000>2\), by its first four positive terms; hence
\(\log2<7/10\). Since \(\sqrt2>7/5\), one has \(\beta<1/2\).
Consequently the complete prime term satisfies \[
\boxed{\int_{-1}^1v(x)|f(x)|^2\,dx-\langle f,B_2f\rangle\ge0.}
\tag{FC19}
\] The elementary estimate \(e<11/4\) used here follows, for example,
from \[
e<\sum_{j=0}^3\frac1{j!}+
\frac{1/4!}{1-1/5}
=\frac{87}{32}<\frac{11}{4}.
\] It also gives \(\log3>1\), verifying the claimed range of prime
terms.

\subsection{5. The global moments give a quantitative Legendre
gap}\label{the-global-moments-give-a-quantitative-legendre-gap}

The nonnegative Hermitian form of \(\mathcal L\) is \[
\mathcal E(f,g)=\frac14\int_{-1}^1\int_{-1}^1
\frac{\overline{f(x)-f(y)}(g(x)-g(y))}{|x-y|}\,dx\,dy.
\] For a monomial, factor \(x^n-y^n\), divide by \(x-y\), and integrate
each term with \(\operatorname{sgn}(x-y)\). This gives \[
\mathcal Lx^n=H_nx^n-
\sum_{j=1}^{\lfloor n/2\rfloor}\frac{x^{n-2j}}{2j},
\qquad H_n=\sum_{j=1}^n\frac1j,\quad H_0=0.
\tag{FC20}
\] Symmetry and degree preservation therefore give
\(\mathcal LP_n=H_nP_n\) for the Legendre polynomials. Put \[
p_n=\sqrt{\frac{2n+1}{2}}P_n,\qquad
c_n=\langle p_n,f\rangle.
\] For a finite Legendre projection \(p^{(N)}=\sum_{n=0}^Nc_np_n\), the
integral form gives \[
\mathcal E(f,f)=
\sum_{n=0}^NH_n|c_n|^2+
\mathcal E(f-p^{(N)},f-p^{(N)}).
\] Indeed \(\mathcal E(p_n,f)=H_n\langle p_n,f\rangle\), by Fubini, and
all these integrals are finite for smooth \(f\) and polynomials.
Nonnegativity and completeness of the Legendre basis imply \[
\mathcal E(f,f)\ge
\frac32\|f\|_2^2-\frac32|c_0|^2-\frac12|c_1|^2.
\tag{FC21}
\] No convergence assumption on an infinite operator expansion is
required.

The two equations (FC12) give orthogonality of \(f_{\rm even}\) to
\(\cosh(ax)\), and of \(f_{\rm odd}\) to \(\sinh(ax)\). With
\(p_0=1/\sqrt2\) and \(p_1=\sqrt{3/2}\,x\), Taylor's integral remainder
and Cauchy--Schwarz give \[
|c_0|^2\le A(T)\|f_{\rm even}\|_2^2,\qquad
|c_1|^2\le\frac5{21}A(T)\|f_{\rm odd}\|_2^2,\qquad
A(T)=\frac{(T/4)^4\cosh^2(T/4)}{20}.
\tag{FC22}
\] The two remainder inequalities are \[
|\cosh(ax)-1|\le\frac{a^2x^2\cosh a}{2},\qquad
|\sinh(ax)/a-x|\le\frac{a^2|x|^3\cosh a}{6}.
\] After squaring, \(\int_{-1}^1x^4dx=2/5\) and \(\int_{-1}^1x^6dx=2/7\)
give the constants \(1/20\) and \(1/84\), respectively.

The constant \(1/4\) in \(r\) gives exactly \[
\langle f,\mathcal R_T^{(0)}f\rangle
=\frac T8\left|\int_{-1}^1f(x)\,dx\right|^2
=\frac T4|c_0|^2.
\tag{FC23}
\] Combining (FC21)--(FC23), and using orthogonality of the even and odd
parts, proves \[
\mathcal E(f,f)-\langle f,\mathcal R_T^{(0)}f\rangle
\ge
\left\{\frac32-\left(\frac32+\frac T4\right)A(T)\right\}\|f\|_2^2.
\tag{FC24}
\] The odd deficit is \(5A(T)/42\), which is bounded by the displayed
even deficit. This retains the constant kernel and uses only the
original two global moments.

\subsection{6. A certified kernel bound with the original
scale}\label{a-certified-kernel-bound-with-the-original-scale}

The analytic function \(r\) at zero has expansion \[
r(z)=\sum_{j=0}^\infty b_jz^j,\qquad
b_j=\frac{2^jB_{j+1}(3/4)}{(j+1)!},
\] \[
(b_0,\ldots,b_8)=
\left(\frac14,-\frac1{48},-\frac1{32},
\frac7{11520},\frac5{1536},-\frac{31}{1935360},
-\frac{61}{184320},\frac{127}{309657600},
\frac{277}{8257536}\right).
\tag{FC25}
\] This follows from \[
\frac{e^{z/2}}{2\sinh z}
=\frac1{2z}\frac{2ze^{3z/2}}{e^{2z}-1}.
\] The nearest nonremovable singularities are \(z=\pm i\pi\), so \(r\)
is analytic on the closed disc of radius \(2\).

For \(z=x+iy\), \(|z|=2\), \[
|\sinh z|^2=\sinh^2x+\sin^2y
\ge x^2+\frac{\sin^22}{4}y^2\ge\sin^22.
\] The first inequality uses the decrease of \(\sin y/y\) on \([0,2]\);
its derivative has the sign of \(y\cos y-\sin y\), whose derivative is
\(-y\sin y<0\). The alternating sine series gives \[
\sin2>2-\frac{2^3}{3!}+\frac{2^5}{5!}-\frac{2^7}{7!}
=\frac{286}{315}>\frac9{10}.
\] Thus, on that circle, \[
|r(z)|\le\frac e{2\sin2}+\frac14
<\frac{16}{9}<2.
\tag{FC26}
\] Cauchy's estimate consequently bounds the degree-eight tail. Define
the explicit rational function \[
E(T)=\sum_{j=1}^8|b_j|T^j+
\frac{2(T/2)^9}{1-T/2}\qquad(0<T<2).
\tag{FC27}
\] Then \(|r(s)-1/4|\le E(T)\) for \(0\le s\le T\). For the nonconstant
remainder operator from (FC13), Schur's row and column bounds give \[
\|\widetilde{\mathcal R}_T\|\le T E(T).
\tag{FC28}
\] Every coefficient and the remainder scale are those of the actual
kernel.

Before substituting a width, (FC14), (FC19), (FC24), and (FC28) prove on
the whole interval (FC16) \[
\boxed{
Q(F)\ge c(T)\|F\|_2^2,\qquad
c(T)=\frac32-\log(\pi T)-\gamma_E
-T E(T)-\left(\frac32+\frac T4\right)A(T).
}
\tag{FC29}
\] This is an explicit bound for every width in a stated numerical
interval whose prime condition has already been proved. It does not
assume the missing sign of a residual operator. If desired, a rational
majorant for the moment quantity is \[
A(T)\le
\frac{(T/4)^4}{20(1-T^2/32)^2},
\] because \((2n)!\ge2^n\) gives \(\cosh(T/4)\le(1-T^2/32)^{-1}\) on
this range.

\subsection{\texorpdfstring{7. Evaluation at
\(T=19/25\)}{7. Evaluation at T=19/25}}\label{evaluation-at-t1925}

Here \(a=19/100\), and \[
\cosh a\le\frac{20000}{19639}<\frac{33}{32}.
\] Exact rational substitution in (FC22) gives \[
A(T_*)<
\frac{141919569}{2048000000000},
\] \[
\left(\frac32+\frac{T_*}{4}\right)
\frac{141919569}{2048000000000}
=\frac{23984407161}{204800000000000}
=\frac1{8000}
-\frac{1615592839}{204800000000000}
<\frac1{8000}.
\tag{FC30}
\] The kernel bound is exactly \[
E(T_*)=
\frac{4199809019840292631}{117180000000000000000}
=\frac9{250}-
\frac{18670980159707369}{117180000000000000000}
<\frac9{250}.
\] Consequently \[
T_* E(T_*)<\frac{171}{6250}.
\tag{FC31}
\]

For the scalar term, the elementary identity \[
\frac{22}{7}-\pi=
\int_0^1\frac{x^4(1-x)^4}{1+x^2}\,dx>0
\] gives \(\pi<22/7\). The exact positive exponential truncation \[
\sum_{j=0}^7\frac{(43/50)^j}{j!}
-\frac{33}{14}
=\frac{1126807544917}{187500000000000}>0
\] proves \(\log(33/14)<43/50\). For Euler's constant, monotonicity of
\(H_n-\log n\), together with the first four terms of (FC18), gives \[
\gamma_E<H_{256}-8\log2
<\frac{612435}{100000}-8\frac{53056}{76545}
=\frac{177361483}{306180000}<\frac{29}{50}.
\tag{FC32}
\] Here \(H_{256}<612435/100000\) is a finite exact rational sum,
checked independently by the accompanying source. Therefore \[
\log(3\pi/4)+\gamma_E<\frac{43}{50}+\frac{29}{50}=\frac{36}{25}.
\] Since \(T_*/(3/4)=76/75\) and \(\log(1+x)<x\) for \(x>0\), \[
\boxed{
\log(\pi T_*)+\gamma_E
<\frac{36}{25}+\frac1{75}=\frac{109}{75}.
}
\tag{FC33}
\]

Substitute (FC30)--(FC33) into (FC29). For nonzero \(F\), \[
\begin{aligned}
Q(F)&>
\left(\frac32-\frac{109}{75}-\frac{171}{6250}-\frac1{8000}\right)
\|F\|_2^2\\
&=\frac{11509}{600000}\|F\|_2^2.
\end{aligned}
\tag{FC34}
\] Its margin above \(1/64\) is exactly \(1067/300000>0\), proving
(FC2).

There is also a separate, slightly weaker rational scalar route. The
five terms \[
\sum_{j=0}^4\frac{(7/8)^j}{j!}
=\frac{78443}{32768}
>\frac{418}{175}>\pi T_*
\] give \(\log(\pi T_*)+\gamma_E<291/200\), and hence \[
Q(F)>\frac{3503}{200000}\|F\|_2^2>
\frac1{64}\|F\|_2^2.
\tag{FC35}
\] The stronger constant in (FC34) exceeds this by exactly \(1/600\).
Both estimates use the same full operator and prime term.

\subsection{8. Exact support range of the fixed first-prime
tests}\label{exact-support-range-of-the-fixed-first-prime-tests}

For the programme's fixed radius \(r=1/64\), the test \(f_r\) has
support in \([-r,r]\). Its autocorrelation has support in \([-2r,2r]\),
so the possible first-prime translation window is \[
\log2-2r\le a\le\log2+2r.
\] The union of the supports of \(f_r\) and \(\tau_af_r\) has containing
interval length at most \(a+2r\). Throughout this window it is therefore
at most \[
\log2+4r=\log2+\frac1{16}.
\tag{FC36}
\] This is below \(19/25\), with a rational proof. The first five
positive terms of \(e^{25/36}\) exceed \(2\), so \(\log2<25/36\). Then
\[
\log2+\frac1{16}<\frac{109}{144}<\frac{19}{25},
\qquad \frac{19}{25}-\frac{109}{144}=\frac{11}{3600}>0.
\tag{FC37}
\] The finite exponential comparison is included in the exact checker.
Thus the larger coercive support interval covers the entire first-prime
window for these original-radius tests, including every linear
combination of the two translates. The complete fixed-test correlation
calculation uses this inclusion; no outside-window bound is inferred
here.

\subsection{9. Original primitives, finite matrices, and supported
zero}\label{original-primitives-finite-matrices-and-supported-zero}

The endpoint-killing operator is unchanged: \[
\mathcal D=\partial_t^2-\frac14,\qquad
M_{\mathcal Dg}(s)=s(s-1)M_g(s).
\tag{FC38}
\] The identity follows by two integrations by parts for
\(g\in C_c^\infty\), and \(\mathcal D\) does not enlarge support.
Conversely every compact smooth \(F\) satisfying the two moments has the
unique compact smooth primitive \[
g(t)=-\int_{\mathbb R}e^{-|t-u|/2}F(u)\,du,\qquad
\mathcal Dg=F.
\tag{FC39}
\] The kernel has derivative jump one. Outside a containing support
interval, its two tails are the original two exponential moments and
hence vanish. A difference of two compact primitives solves the
homogeneous equation and must be zero. Thus (FC39) preserves the
containing support interval.

The norm calculation gives \[
\|\mathcal Dg\|_2^2=
\|g''\|_2^2+\frac12\|g'\|_2^2+\frac1{16}\|g\|_2^2.
\] Consequently, for nonzero \(g\) with support diameter at most
\(19/25\), \[
Q(\mathcal Dg)>
c_*\left(\|g''\|_2^2+\frac12\|g'\|_2^2+\frac1{16}\|g\|_2^2\right).
\tag{FC40}
\] The interval coordinate map has the exact intertwining formula \[
U_{T,c}\mathcal D U_{T,c}^{-1}
=\frac4{T^2}\partial_x^2-\frac14.
\tag{FC41}
\] No constant has been absorbed into a different differential operator.

For a finite family in one common containing interval of length
\(19/25\), the full sesquilinear form satisfies \[
[Q(F_i,F_j)]\succeq c_*[\langle F_i,F_j\rangle].
\tag{FC42}
\] Apply (FC34) to \(\sum c_jF_j\) to prove the matrix statement; the
difference is strictly positive whenever that function is nonzero.

Finally retain the entire original support space \(\mathbb C[L]\). With
\(H=M_{F^\#*F}\), its distributions are \[
\begin{aligned}
\boldsymbol B_L&=(H(0)+H(1))\mathbf e_{1_L}
+H(0)\sum_{\lambda\ne1_L}\mathbf e_\lambda,\\
\boldsymbol Z_L&=Q(F)\mathbf e_{1_L},\\
\boldsymbol D_L&=(P_{\rm fin}-A_\infty)\mathbf e_{1_L}
+H(0)\sum_{\lambda\ne1_L}\mathbf e_\lambda,\\
\boldsymbol B_L-\boldsymbol Z_L&=\boldsymbol D_L.
\end{aligned}
\tag{FC43}
\] Both endpoint coefficients vanish by (FC5); their coordinate spaces
remain. Every linear map displayed above has the fixed-label lift
\((F,\lambda)\mapsto(TF,\lambda)\) on its stated domain. It commutes
with amplitude projection and leaves the support label unchanged. A zero
amplitude in such a carrier remains its supported zero \((0,\lambda)\);
it is not identified with the unsupported element \(\tau\).

\subsection{Source and verification
record}\label{source-and-verification-record}

The starting \(3/4\) theorem and interval-operator method were supplied
in the companion manuscript ``Direct RH continuation: the optimal packet
cost in the full Weil form,'' N1--N11 and N48--N52, dated 23 September
2026, source SHA256
D3B1AC08C71F6104CFE011BA5E5B86F6FBE182916C0B0270511671B1F7165D5B. The
independent audit is retained separately. The present proof gives all
arguments needed for the enlarged interval and its stronger constant.

The programme's original full formula is SZW19--SZW38, with its human
Weil and Connes sources; HA13 supplies the real kernel (FC7), and
PS1--PS15 supplies the original support and primitive conventions, which
are rederived where used here. No historical priority claim is made for
the interval kernel, Legendre diagonalization, or local Weil positivity.

The accompanying exact-rational checker divides the kernel power series
independently, verifies every displayed rational estimate at \(3/4\) and
\(19/25\), evaluates the finite harmonic sum, and checks the exponential
comparisons. It uses rational arithmetic only. These checks supplement
the analytic estimates (FC7)--(FC43); they do not replace the
convergence, domain, or prime-interaction proofs.

\end{document}
